\begin{Theorem}
\label{thm:sobolev-banach-separability}
空间 $W^{k,p}(\Omega)$ 为 Banach 空间. 若 $p<\infty$, 则它可分.
\end{Theorem}

\begingroup
\renewcommand{\endProof}{\par\nobreak\hfill$\blacksquare$\par\vspace{0.1cm}}
\begin{Proof}
对 $u\in W^{k,p}(\Omega)$, 令 $u_\alpha=D^\alpha u$. 那么映射 $u\to(u_\alpha)_{|\alpha|\leq k}$ 是从 $W^{k,p}(\Omega)$ 到直积 $\prod_{|\alpha|\leq k}L^p(\Omega)$ 的一个子空间 $M$ 上的等距映射. 只要证明 $M$ 闭, 就得到 $W^{k,p}(\Omega)$ 的完备性. 为此需要证明, 若 $u_n\in W^{k,p}$ 且 $D^\alpha u_n\to v^\alpha$ 在 $L^p(\Omega)$ 中成立, 则 $v^\alpha=D^\alpha v^0$. 但这显然成立, 因为 $L^p$ 收敛蕴含分布收敛.

若 $p<\infty$, 则 $\prod_{|\alpha|\leq k}L^p(\Omega)$ 可分. 因而只需证明可分 Banach 空间的子空间可分. 以下引理证明这一点, 从而完成本定理的证明.
\end{Proof}
\endgroup

\begin{Lemma}
\label{lem:separable-banach-linear-subspace}
设 $X$ 为可分 Banach 空间, $M$ 为 $X$ 的线性子空间. 那么 $M$ 可分.
\end{Lemma}

\begin{Proof}
对任意给定的 $\varepsilon>0$, $X$ 中存在可数集 $\{x_n\}_{n\in\mathbb{N}}$, 使对任意 $x\in X$, 都存在 $n\in\mathbb{N}$, 满足 $\|x-x_n\|\leq\varepsilon$. 从序列 $x_n$ 中删去所有使以 $x_n$ 为中心, 半径为 $\varepsilon$ 的球与 $M$ 不相交的元素. 剩余点组成新序列 $y_n$. 现在设 $u_n$ 为 $M$ 中位于 $y_n$ 的 $\varepsilon$ 邻域内的任意点. 那么 $M$ 中每个点与某个 $y_n$ 的距离不超过 $\varepsilon$, 因而与某个 $u_n$ 的距离不超过 $2\varepsilon$. 让 $\varepsilon$ 遍历任意收敛到零的序列, 就可生成 $M$ 中的一个可数稠密集.
\end{Proof}

若 $p<\infty$, 则 $\mathcal{D}(\Omega)$ 在 $L^p(\Omega)$ 中稠密. 对空间 $W^{k,p}(\Omega)$, 一般并非如此. 事实上有以下结果.

\begin{Lemma}
\label{lem:test-functions-sobolev-nondensity}
设 $\Omega$ 有界, $k\geq1$ 且 $1\leq p\leq\infty$. 那么 $\mathcal{D}(\Omega)$ 在 $W^{k,p}(\Omega)$ 中不稠密.
\end{Lemma}

\begin{Proof}
若 $u$ 为 $C^{\infty}(\overline\Omega)$ 中的任意函数, 则映射
\begin{equation}\label{eq:sobolev-test-functions-annihilating-functional}
L:\phi\mapsto\int_{\Omega}\sum_{|\alpha|\leq k}D^\alpha\bar u(\bm{x})D^\alpha\phi(\bm{x})\,\mathrm{d}\bm{x}
\end{equation}
定义了 $W^{k,p}(\Omega)$ 上的有界线性泛函. 此外, 若
\begin{equation}\label{eq:sobolev-nondensity-annihilating-elliptic-equation}
\sum_{|\alpha|\leq k}(-1)^{|\alpha|}D^{2\alpha}\bar u=0,
\end{equation}
则对每个 $\phi\in\mathcal{D}(\Omega)$ 都有 $L(\phi)$ 为零. 在包含 $\Omega$ 的任意平行多面体中, 容易用分离变量法找到 \eqref{eq:sobolev-nondensity-annihilating-elliptic-equation} 的一个非零解 $u$; 见 Chapter~\ref{chap:introduction}. 此外, $L(u)>0$ 且 $u\in W^{k,p}(\Omega)$; 因而 $\mathcal{D}(\Omega)$ 不可能稠密.
\end{Proof}

上一个引理引出以下定义.

\begin{Definition}
\label{def:sobolev-test-function-closure}
用 $W_0^{k,p}(\Omega)$ 表示 $\mathcal{D}(\Omega)$ 在 $W^{k,p}(\Omega)$ 中的闭包.
\end{Definition}
